\section{What the criticism scores measure}
\label{app:metrics}
This appendix distinguishes three questions relevant to interpreting the scores in Section~\ref{sec:evaluation}. They are whether a formal counterexample refutes the source claim, whether a critic selected a false assertion, and whether the critic's explanation is valid. The propositions are elementary, and we developed their written proofs with AI tools; they are separate from the Lean-certified mathematics elsewhere in the paper. The worked comparison uses the saved reviews of 24 texts from an earlier evaluation set.

\subsection{The logical connection from a formal witness to a source claim}
\label{app:transport}
Section~\ref{sec:evidence} defines the source assertion $Q_s$, the formal witness, and its interpretation $e$, and Equations~\eqref{eq:source-claim}--\eqref{eq:transport} state the source-correspondence requirements. The proposition below is the implication used there.

\begin{proposition}[One-sided refutation transport]
\label{prop:transport}
If Equation~\eqref{eq:formal-witness} and the witness-level implications in Equation~\eqref{eq:transport} hold, then $A_s(e(w))\wedge\neg C_s(e(w))$, so the source assertion $Q_s$ is false.
\end{proposition}
\begin{proof}
The first implication gives the source assumptions. If the source conclusion were true, the second implication would make the encoded conclusion true, contradicting the certificate. Hence the source conclusion is false at an allowed source object.
\end{proof}
Full equivalence between the formal and source statements is one way to satisfy these conditions, not the only one. The sorting example in Section~\ref{app:source-admissibility} illustrates the first condition: the witness must satisfy the source's endpoint inequality. Writing down $X_s,A_s,C_s$, and $e$ makes the interpretation a reviewer must judge explicit.

\subsection{Secondary reviews of generated objections}
\label{app:secondary-objections}
The Sonnet miss analysis in Section~\ref{app:ccl_miss_diagnostic} shows that a response can miss the designated reference line and still identify another false assertion. The study below checks, for each saved objection, whether it identifies its target and whether its evidence is valid.

\paragraph{Experiment B: LLM review of saved responses from an earlier evaluation set.} Experiment~B examined 48 saved responses, one from Gemini~3.1 Pro and one from Nemotron~3 Super 120B for each of 24 sampled texts of the earlier evaluation set. Two LLM reviewers separately read anonymized packets and judged whether each explanation challenged the selected claim (J1) and whether its mathematical evidence was valid (J3).

The question is whether a response both identifies the assertion it challenges and supplies valid evidence under that assertion's assumptions. Table~\ref{tab:b-results} reports this joint endpoint after the recorded LLM adjudications and the checks that each objection applies to the selected source claim under its assumptions; the first-reader judgments designated in advance are kept below. Six of the 10 reference misses supplied valid objections under the joint criterion.

\begin{table}[h]
\centering
\begin{tabular}{l c c c c}
\toprule
\textbf{Model} & \textbf{Reference} & \textbf{Valid obj.} & \textbf{Valid obj.} & \textbf{Total valid} \\
& \textbf{matches} & \textbf{among matches} & \textbf{among misses} & \textbf{objections} \\
\midrule
Gemini 3.1 Pro & 20 & 18 & 2 & 20 / 24 \\
Nemotron 3 Super 120B & 18 & 15 & 4 & 19 / 24 \\
\bottomrule
\end{tabular}
\caption{LLM review of joint target identification and evidence validity on 24 paired texts. A counted valid objection is a response that met the rules, challenges its selected claim, and supplies mathematically valid evidence under the source assumptions, after the recorded adjudications and the checks that the objection applies to the selected source claim. Table~\ref{tab:b-results-appendix} gives the first-reader judgments and the counts across the recorded review stages.}
\label{tab:b-results}
\end{table}

\paragraph{Detailed judgments.} The upper panel of Table~\ref{tab:b-results-appendix} gives the first-reader evidence-validity judgments (J3) designated in advance; the lower panel shows the joint positive-evidence counts across the recorded review stages: the first designated LLM reviewer alone, agreement of both reviewers, inclusion of the recorded adjudications, and finally the source-assumption checks.

\begin{table}[h]
\centering\small
\begin{tabular}{l l r r r}
\toprule
\multicolumn{5}{l}{\emph{First-reader evidence-validity judgments (J3) for the 48 saved responses}}\\
\textbf{Model} & \textbf{First-choice outcome} & \textbf{Valid} & \textbf{Invalid} & \textbf{Unresolved} \\
\midrule
Gemini 3.1 Pro & Hit@1 & 18 & 0 & 2 \\
Gemini 3.1 Pro & Miss & 3 & 1 & 0 \\
Nemotron 3 Super 120B & Hit@1 & 16 & 2 & 0 \\
Nemotron 3 Super 120B & Miss & 5 & 0 & 1 \\
\bottomrule
\end{tabular}

\smallskip
\begin{tabular}{l c c}
\toprule
\multicolumn{3}{l}{\emph{Responses meeting the rules that satisfy both J1 and J3}}\\
\textbf{Review stage} & \textbf{Gemini 3.1 Pro} & \textbf{Nemotron 3 Super 120B} \\
\midrule
First reader & 18 / 24 & 19 / 24 \\
Both readers & 18 / 24 & 18 / 24 \\
After LLM adjudication & 20 / 24 & 20 / 24 \\
With source-assumption check & 20 / 24 & 19 / 24 \\
\bottomrule
\end{tabular}
\caption{Experiment B in detail. Upper panel: the first-reader evidence-validity judgments (J3), by whether the response's first choice was a reference line. Lower panel: the number of responses meeting the rules that satisfy both target identification (J1) and evidence validity (J3) across the recorded review stages (ending with the final joint endpoint after the source-assumption check).}
\label{tab:b-results-appendix}
\end{table}

A third LLM reviewer examined eight flagged responses. P06 contained a complete explanation, although the attempt exceeded the allowed budget. P43 supplied valid mathematical evidence but cited the wrong line. P35 challenged a stronger condition than the selected statement claimed, so its counterexample did not address that claim. For P31, the calculation refuted the implication stated in the model's explanation, while the supplied example violated two assumptions in the source. Applying the study's rule that counterexamples must satisfy the source assumptions classifies P31's evidence as invalid.

\subsection{Partial labels and generated-objection soundness}
A reference-selection score measures agreement with the designated target. Two broader outcomes are finding any false assertion and supplying a valid counterexample for it. Because the reference labels cover selected assertions only and the score does not assess explanations, the available data bound these outcomes without determining them, a partial identification problem~\citep{manski2003partial}.

In this subsection $N$ is the number of benchmark inputs and all sums over $i$ run from one to $N$. For item $i$, let $F_i$ be the set of admissible line labels that locate genuinely false assertions under the source assumptions (a label may identify part of a displayed proposition); this set describes the truth of the assertions and need not be fully annotated. Let $G_i$ be the set of reference lines and $V_i$ the indicator that the response met the rules. For a first-choice label $\ell_i$, define
\[
h=\frac1N\sum_i V_i\ind\{\ell_i\in G_i\},\quad
q=\frac1N\sum_i V_i\ind\{\ell_i\in F_i\},\quad
c=\frac1N\sum_i V_i.
\]
The fractions $h$ and $q$ count responses that meet the rules and select, respectively, a reference line or any genuinely false assertion; $c$ counts all responses that meet the rules. All three use the full $N$-input denominator. Unlike $h$ and $c$, $q$ is not observable in general, because the full sets $F_i$ are not annotated. Null labels belong to neither set.
\begin{proposition}[Partial identification under sound references]
\label{prop:partial}
If $G_i\subseteq F_i$ for every item, then $h\le q\le c$. If the inclusion may fail on at most $m$ items, then
\[
\max(0,h-m/N)\le q\le c.
\]
Let $E_i$ indicate that the critic's supplied objection is valid for its selected target, and define $r=N^{-1}\sum_i V_i\ind\{\ell_i\in F_i\}E_i$. Without evidence about $E_i$, no strictly positive lower bound on $r$ follows from $h$ alone.
\end{proposition}
\begin{proof}
On a sound item, reference membership implies false-assertion membership. At most $m$ reference hits can fail that implication on the remaining items. No valid target can be credited without a valid output, giving the upper bound $c$. Finally, every $E_i$ can be zero even when every chosen target matches its reference set: a critic may attach an invalid witness to the correct line. Hence $r=0$ is compatible with positive or perfect reference localization.
\end{proof}
The same argument applies to top-$k$ selection, with explanation validity assessed separately. When the designated references are sound, matching one is evidence that the critic selected a false assertion, and a miss can still identify a different false assertion, as the diagnostic in Section~\ref{app:ccl_miss_diagnostic} shows; estimating $q$ or $r$ requires the corresponding target and explanation judgments.

\subsection{Comparing critics with partially reviewed line selections}
\label{app:sharp-comparison}
If two critics select the same line, its truth status affects both scores equally, so only the selections on which they disagree can change their score difference. The following result turns that observation into exact bounds.

Set a critic's first-choice label to $\bot$ when its response does not meet the rules or its first-choice label is invalid, so that critic $a$ or $b$ supplies $p_i^a,p_i^b\in L_i\cup\{\bot\}$, where $L_i$ is the addressable line set. Write $z_{i\ell}=\ind\{\ell\in F_i\}$ and $z_{i\bot}=0$. A line reviewed as false under the source assumptions has value one; a line reviewed as not false has value zero. Every other indicator is unknown. A reference label has value one only when its validity under the source assumptions is assumed or established.

For critic $m\in\{a,b\}$, let $\underline z_i^m,\overline z_i^m$ be the lower and upper possible values of $z_{i,p_i^m}$. They are $(1,1)$ for known false assertions, $(0,0)$ for known nonfalse or null predictions, and $(0,1)$ for unadjudicated lines. Indicators of the same line on the same item are shared, not assigned separately to each critic.

\begin{proposition}[Sharp comparison in the line-label model]
\label{prop:sharp-comparison}
Over assignments of unknown line indicators consistent with these known labels and with no further constraints, the false-assertion localization difference
\[
\Delta_F=\frac1N\sum_i\bigl(z_{i,p_i^a}-z_{i,p_i^b}\bigr)
\]
has sharp lower and upper bounds
\begin{align}
L&=\frac1N\sum_{i:p_i^a\ne p_i^b}(\underline z_i^a-\overline z_i^b),&
U&=\frac1N\sum_{i:p_i^a\ne p_i^b}(\overline z_i^a-\underline z_i^b).
\label{eq:sharp-comparison}
\end{align}
Both endpoints are attained in this model. ``Sharp'' means that no narrower interval follows from these labels alone under the stated absence of further constraints.
\end{proposition}
\begin{proof}
Coincident predictions share an indicator, so their contribution is exactly zero even when it is unknown. For distinct predictions the difference is at least $\underline z_i^a-\overline z_i^b$ and at most $\overline z_i^a-\underline z_i^b$. To attain the lower bound, set each unknown indicator selected by $a$ to zero and each selected by $b$ to one on disagreement items; reverse these assignments for the upper bound. Distinct selected lines impose no conflicting assignment within an item, and this model imposes no cross-item constraints. Thus the extremal choices can be made simultaneously.
\end{proof}

\paragraph{Which additional reviews narrow the comparison.}
The identification width is
\[
U-L=\frac1N\sum_{i:p_i^a\ne p_i^b}
\bigl[(\overline z_i^a-\underline z_i^a)+(\overline z_i^b-\underline z_i^b)\bigr].
\]
Only unknown predictions on disagreement items contribute: one unknown line gives width $1/N$, two give $2/N$, and a shared unknown prediction gives zero. Reviewing one such line narrows the interval whichever label it receives, so the disagreements identify the additional source judgments that matter most for the paired comparison.

The bounds are exact for the stated model, which permits every assignment consistent with the known line labels; logical relationships among assertions or repeated sources can narrow them further. An identification interval describes the possible score differences while some labels are unknown, whereas a sampling confidence interval describes uncertainty from sampling texts, and reviewing explanation validity answers a third question, whether the critic supplies valid evidence for its chosen assertion. Applying the bounds requires item-level predictions and judgments.

\paragraph{Numerical application on 24 texts.} The proposition concerns which selected assertions are false, not whether the critic's own explanation proves them false. Using the saved judgments and source checks, we assign one shared truth indicator to each $(\text{input}, \text{line})$ pair: a valid objection can establish a positive line label, while an invalid explanation alone does not establish a negative one.

Gemini and Nemotron make the same normalized prediction on 15 of 24 inputs; their contribution is exactly zero regardless of the line's status. They disagree on nine inputs. Let $z_{19}$ and $z_{43}$ be the unknown false-assertion indicators for Nemotron's P19 line 24 and P43 line 230. The seven other disagreement contributions sum to zero based on the accepted labels or null predictions. These two contribute $-z_{19}$ and $1 - z_{43}$. Thus the Gemini-minus-Nemotron score difference is $(1 - z_{19} - z_{43})/24$, between $-1/24$ and $1/24$.

Positive line labels are accepted conditionally from the saved LLM judgments, all other non-null labels stay unknown, the references are not assumed sound, and no negative label is inferred from invalid evidence. These are the 24 paired texts of Experiment B, separate from the current 100-input benchmark. The interval $[-1/24,1/24]$ is conditional on the recorded LLM judgments and summarizes the two unknown labels; it is not a 95\% confidence interval. The saved inputs, judgments, and reproducible calculation are documented in \path{studies/criticism_endpoints/README.md}.

\subsection{Sensitivity to revised reference lines}
Let $k$ be the number of candidate slots scored. Fix $P_i^k$, the set of valid line labels in the original first $k$ slots for input $i$; null slots remain unfilled and later predictions are not promoted. Write $G_i\triangle G_i' $ for labels that belong to exactly one of the old and new reference sets. Let $G_i,G'_i$ be old and revised reference sets and $r=|\{i:G_i\ne G'_i\}|$.
\begin{proposition}[Prediction-aware reference perturbation]
\label{prop:perturb}
For fixed outputs and a fixed $N$-item population,
\begin{align}
\left|\Hit_G@k-\Hit_{G'}@k\right|
&\le \rho(k):=\frac1N\sum_i
\ind\{P_i^k\cap(G_i\triangle G'_i)\ne\varnothing\}
\le\frac rN.
\label{eq:prediction-aware-revision}
\end{align}
For two critics $a,b$, their paired score difference changes by at most $\rho_a(k)+\rho_b(k)\le2r/N$.
\end{proposition}
\begin{proof}
An item cannot change hit status unless a selected label lies in the symmetric difference $G_i\triangle G'_i$. Even then, a retained reference hit may keep its status unchanged. Each binary hit changes by at most one, so summing gives the first inequality. Only items whose reference set changed can contribute to $\rho$, giving the second. Applying the triangle inequality to two critics yields the paired bound. The coarse constants are attainable: a changed reference can turn one miss into a hit, or exchange reference membership between two distinct first-choice predictions, changing their difference by two.
\end{proof}
The prediction-aware bound counts outcomes that could change, not net changes: a gained hit and a lost hit can cancel in an aggregate score, and an edit away from every selected line cannot affect the score at all. Counting changed reference sets alone can therefore overstate the possible effect of a revision.

The bound applies to any recorded revision of the checked line sets, and rescoring the saved predictions against both versions gives the actual change. If $m$ excluded inputs are misses for every reporting view, each rate changes from $H/N$ to $H/(N-m)$ while the hit count $H$ stays fixed. The release provides the line sets, exclusion lists, and per-input predictions for these calculations.

\subsection{Why paired comparisons require individual outcomes}
Suppose two models have $a$ and $b$ hits on the same $N$ items, with $a\ge b$. Their common-hit count $t$ can lie anywhere in
\[
\max(0,a+b-N)\le t\le b.
\]
The discordant counts are $a-t$ and $b-t$, so a paired test needs more than the marginals. In the current 100-input results, Gemini has 73 first-choice matches and Opus 68; these totals allow $t$ anywhere from 41 to 68, that is, between 5 and 59 inputs on which exactly one model succeeds. The released per-input predictions determine the actual overlap. The main text therefore reports the highest observed scores, and paired comparisons use the item-level outcomes.

\subsection{Costs with incomplete usage records}
Let the total nonnegative costs in two conditions be $C_A=L_A+U_A$ and $C_B=L_B+U_B$, where $L_A,L_B$ are recorded subtotals and $U_A,U_B\ge0$ unrecorded usage. For positive denominators, $L_A/L_B$ bounds $C_A/C_B$ in neither direction in general, because either unrecorded remainder can be the larger one. A budget charge is an accounting debit, not a measurement of that remainder, and a cost total over selected successful attempts describes a different set of executions from one that includes failures and retries. We therefore name the included calls and report measured usage, lower bounds, and budget charges separately.
